%% Copernicus Publications Manuscript Preparation Template for LaTeX Submissions
%% ---------------------------------
%% This template should be used for copernicus.cls
%% The class file and some style files are bundled in the Copernicus Latex Package, which can be downloaded from the different journal webpages.
%% For further assistance please contact Copernicus Publications at: production@copernicus.org
%% https://publications.copernicus.org/for_authors/manuscript_preparation.html


%% Please use the following documentclass and journal abbreviations for preprints and final revised papers.

%% 2-column papers and preprints
\documentclass[hess, manuscript]{copernicus}


%% \usepackage commands included in the copernicus.cls:
%\usepackage[german, english]{babel}
\usepackage{tabularx}
%\usepackage{cancel}
%\usepackage{multirow}
%\usepackage{supertabular}
%\usepackage{algorithmic}
%\usepackage{algorithm}
%\usepackage{amsthm}
%\usepackage{float}
%\usepackage{subfig}
%\usepackage{rotating}
\usepackage{pdflscape}
\usepackage[section]{placeins}
\usepackage{booktabs}
\usepackage{threeparttable}
\usepackage[dvipsnames]{xcolor}
\usepackage{makecell}
\usepackage{siunitx}

\begin{document}

\title{An algorithmic framework for estimating snow slab mechanical properties from field observations}

% \Author[affil]{given_name}{surname}

\Author[1]{Mary Kate}{Connelly}
\Author[2]{Philipp}{Rosendahl}
\Author[3]{Valentin}{Adam}
\Author[1]{Samuel}{Verplanck}

\affil[1]{205 Cobleigh Hall, Montana State University, Bozeman, MT, USA 59717}
\affil[2]{TU Darmstadt} %% TODO: full address needed
\affil[3]{SLF Davos}    %% TODO: full address needed

\correspondence{Samuel Verplanck (samuelverplanck@montana.edu), Mary Kate Connelly (connellymarykate@gmail.com)}

\runningtitle{Algorithmic framework for snow mechanical properties}

\runningauthor{Connelly et al.}


\begin{abstract}

% STRUCTURE: finding first, then method.
%
% Opening (the problem and why it matters, ~2 sentences):
%   Mechanical stability models for snow slab avalanche release require elastic slab
%   properties -- in particular the bending stiffness D11 -- that cannot be measured
%   directly in the field and must instead be estimated by chaining published
%   parameterizations through standard snowpit observations. How reliably these
%   properties can be estimated, and from what fraction of real field datasets, has
%   not been systematically quantified.
%
% Headline results (lead with the three numbers, ~3 sentences):
%   We apply a systematic framework to 14,951 ECTP slab stratigraphies from
%   10 years of SnowPilot data (51,138 pits, 1 September 2014 -- 31 August 2025)
%   and find: (1) the highest-coverage calculation pathway produces a valid D11
%   estimate for only ~5% of slabs, because the field observations required by any
%   single pathway are rarely all present; (2) for slabs where multiple pathways
%   succeed, D11 varies by a factor of 3--4 depending entirely on method choice,
%   with E-modulus parameterization as the dominant contributor; and
%   (3) within-pathway relative uncertainty from input measurement propagation
%   alone is 40--90%, and compounds when parameterizations are chained.
%
% Implication (the risk argument, ~2 sentences):
%   Together, these results quantify a fundamental constraint on field-data-driven
%   mechanical models: the inputs required by these models are both rarely achievable
%   from standard observations and, when achievable, highly sensitive to
%   parameterization choice. Conclusions about slope stability drawn from mechanical
%   model outputs should account for this dual source of uncertainty.
%
% Method/tool (supporting, ~1--2 sentences):
%   To enable this analysis, we present SnowPyt-MechParams, an open-source Python
%   library that encodes parameterizations as a directed acyclic graph (DAG) and
%   automatically enumerates all 32 valid D11 calculation pathways, propagating
%   input measurement uncertainty through each via first-order Taylor expansion.
%   The library is openly available and designed to incorporate new parameterizations
%   as they are published.

\end{abstract}


\section{Introduction}

\subsection{Mechanical models of snow stability and their input parameters}

% STRUCTURE: open with the operational/risk problem, then trace the chain back to
% the fundamental gap. The framework is introduced later as the method; this section
% establishes what is at stake.
%
% Paragraph 1 -- What mechanical models do and why D11 matters:
%   Introduce the mechanics of slab avalanche release and the role of mechanical models.
%   Cover: crack initiation in the weak layer, crack propagation, and the elastic slab
%   properties that govern the energy available to drive propagation.
%   Identify D11 (bending stiffness) as a key target parameter:
%     - D11 governs the elastic energy stored in the slab and available to drive crack
%       propagation through the weak layer
%     - D11 is a required input to energy-based and combined energy--stress stability
%       models (e.g., Heierli et al. 2008, Rosendahl & Weissgraeber 2020, Siron et al. 2023)
%     - D11 is computed from layer-level elastic properties (Young's modulus E and
%       Poisson's ratio nu) via classical laminate theory
%
% Paragraph 2 -- The estimation problem:
%   E and nu are seldom measured directly in field conditions; they must be estimated
%   by chaining parameterizations from directly observable quantities (hand hardness,
%   grain form, grain size, directly measured density). Density is a shared parameter
%   feeding both E and nu estimation. Chaining introduces two compounding problems:
%   (a) each link has its own measurement uncertainty, which propagates and amplifies;
%   (b) multiple published parameterizations exist for each step, each with different
%   input requirements, applicable grain forms, and density ranges -- so the result
%   depends on which methods are chosen.
%
% Paragraph 3 -- The gap and its risk implications (state explicitly):
%   No systematic analysis has quantified (a) what fraction of real field observations
%   contain sufficient data to complete any calculation pathway to D11, or (b) how much
%   D11 varies across valid pathways applied to the same field data. Without this,
%   practitioners using field-data-driven mechanical models have no basis for assessing
%   whether their inputs are representative or how sensitive their stability conclusions
%   are to parameterization choice. This paper provides that quantification at population
%   scale for the first time.


\subsection{Snowpits and the SnowPilot database}
\label{sec:snowpilot_database}

% Describe what a snowpit is and what is recorded: layering, grain forms, stability
% tests, terrain and weather conditions. Reference standard observation protocols

A snowpit is a trench excavated into the snowpack, typically to the ground or to a depth of interest. For each stratigraphic layer encountered, the observer records hand hardness, grain form, grain size, and layer thickness; density may also be measured directly on individual layers using a density cutter and scale. Terrain characteristics (slope angle, aspect, elevation) and weather conditions at the time of the observation supplement the stratigraphy. Stability tests are commonly conducted alongside the snow profile -- most frequently the Compression Test (CT), Extended Column Test (ECT), and Propagation Saw Test (PST) -- and record the stress required to initiate a fracture, whether that fracture propagates, and the depth and character of the failure plane. Standard observation protocols are described in \citet{fierz2009}, \citet{greene2022}, \citet{caatechnicalcommittee2024}, and \citet{norwegianwaterresourcesandenergydirectorate2022}.

% Introduce the SnowPilot database (chabot2004, chabot2016):
%   - Free, web-based application for recording and storing snowpit data
%   - Data stored in CAAML XML format (caatechnicalcommittee2024)
%   - Community-sourced: wide range of observer expertise and completeness

SnowPilot \citep{chabot2004,chabot2016} is a free, web-based application for recording, visualising, and storing snowpit observations. Data are archived in the Canadian Avalanche Association Markup Language (CAAML) XML format \citep{caatechnicalcommittee2024}, an open community standard for snow and avalanche data exchange. The database is community-sourced, with contributors spanning avalanche centre forecasters, ski patrol, field researchers, and backcountry educators, and is particularly widely used in North America.

% Present Table (dataset summary) covering 1 September 2014 to 31 August 2025.

A summary of the SnowPilot data used in this analysis -- the decade from 1~September~2014 to 31~August~2025 -- is provided in Table~\ref{tab:snowpilot}.

\begin{table}
\begin{threeparttable}
\caption{Summary of the SnowPilot database for the analysis period 1~September~2014 to 31~August~2025.}
\label{tab:snowpilot}
\begin{tabular}{lr}
\toprule
 & Count \\
\midrule
\multicolumn{2}{l}{\textit{Snowpits}} \\
\quad Total & 51{,}138\tnote{a} \\
\quad Unique observers & 5{,}440 \\
\quad Countries represented & 35 \\
\quad With slope angle observation & 46{,}226 (90.4\,\%) \\
\quad Recorded near an avalanche & 1{,}586 (3.1\,\%) \\
\addlinespace
\multicolumn{2}{l}{\textit{Snow layers}} \\
\quad Total & 377{,}729 \\
\quad With hand hardness & 342{,}690 (90.7\,\%) \\
\quad With primary grain form & 309{,}232 (81.9\,\%) \\
\quad With grain size & 179{,}342 (47.5\,\%) \\
\quad With direct density measurement\tnote{b} & 10{,}692 (2.8\,\%) \\
\addlinespace
\multicolumn{2}{l}{\textit{Stability tests}} \\
\quad Compression tests (CT) & 52{,}305 \\
\quad Extended column tests (ECT) & 48{,}391 \\
\quad Propagation saw tests (PST) & 6{,}326 \\
\bottomrule
\end{tabular}
\begin{tablenotes}
\footnotesize
\item[a] One file (\texttt{snowpits-55240-caaml.xml}) failed to parse and was excluded, leaving 51{,}138 snowpits.
\item[b] SnowPilot allows density to be entered either as an attribute of a stratigraphic layer or as a separate density profile. Our analysis uses only density observations whose depth and thickness match a stratigraphic layer boundary; the count reported here reflects this aligned subset.
\end{tablenotes}
\end{threeparttable}
\end{table}

% Close with a bridging sentence: SnowPilot's scale (51,138 pits, 10 years) makes it
% uniquely suited for population-level quantification of coverage and parameterization
% sensitivity -- the questions this paper addresses. Its community-sourced nature also
% means it reflects the realistic distribution of field observation completeness, rather
% than data collected under controlled research conditions.

The scale of the SnowPilot dataset -- 51{,}138 snowpits over ten years, spanning 35 countries and 5{,}440 unique observers -- makes it uniquely suited for population-level quantification of observational coverage and parameterization sensitivity. Critically, because the data are community-sourced rather than collected under controlled research conditions, coverage fractions derived from this dataset reflect the realistic distribution of field observation completeness encountered by practitioners, rather than the more complete data typical of targeted research campaigns.

\FloatBarrier

\subsection{Objectives and scope}

% STRUCTURE: state the finding-oriented objectives first; the framework is presented
% as the method that enables them, not as an end in itself. SnowPylot is prior work
% (cited in §2.1), not a contribution of this paper.
%
% This paper has two objectives:
%
%   1. FINDING (primary objective): Quantify, at population scale, (a) the fraction of
%      real snowpit observations that contain sufficient data to estimate D11 via any
%      published parameterization pathway, and (b) the sensitivity of the resulting
%      D11 estimates to parameterization choice and to input measurement uncertainty.
%      Applied to 14,951 ECTP slab stratigraphies from 10 years of SnowPilot data,
%      we show that both coverage and consistency are far lower than has been
%      previously recognized.
%
%   2. TOOL (enabling objective): Present SnowPyt-MechParams, an open-source Python
%      library implementing a parameterization graph and execution engine that
%      systematically enumerates and executes all valid calculation pathways from
%      snowpit observations to target mechanical parameters, propagating input
%      measurement uncertainty through each pathway. The library is openly available
%      and designed to support future extensions as new parameterizations are published.
%
% Scope:
%   - Included: slab layer parameters (density, E, nu, shear modulus) and slab-level
%     plate theory parameters (A11, B11, D11, A55)
%   - Excluded: weak layer parameters (shear strength, fracture toughness) -- out of scope
%   - No new parameterizations are developed; existing published methods are implemented
%     and compared
%   - No claim is made about the relationship between estimated D11 and ECT outcome;
%     ECTP slabs serve as a source of field-observed slab stratigraphies
%
% Paper structure:
%   Sect. 2 describes the data, uncertainty representation, parameterizations,
%   the computational framework, and the ECT slab construction procedure.
%   Sect. 3 presents results for density, E, nu, and D11, with coverage and
%   uncertainty quantified at each level.
%   Sect. 4 discusses what these results mean for field-data-driven mechanical models
%   and for the practice of stability assessment.
%   Sect. 5 presents conclusions.


\section{Data and methods}
\label{sec:data_methods}

\subsection{Data: SnowPilot and SnowPylot}
\label{sec:snowpylot}

% Introduce SnowPylot (connelly2025) as an open-source Python library
% that parses SnowPilot CAAML XML files into
% structured Python objects (SnowPit, snow_profile, stability_tests, etc.).

SnowPylot \citep{connelly2025} is an open-source Python library for parsing SnowPilot CAAML XML files into structured Python objects. It provides the data ingestion layer for this analysis, translating the raw CAAML archive into an accessible object hierarchy from which measured field observations are extracted for parameterization.

% Briefly describe the object hierarchy relevant to this analysis:
%   SnowPit -> snow_profile -> layers (hand_hardness, grain_form, thickness, grain_size)
%   SnowPit -> stability_tests -> ECT results (score, failure_depth)

Each CAAML file is parsed into a \texttt{SnowPit} object with two main sub-objects. The \texttt{snow\_profile} object contains the ordered list of stratigraphic \texttt{Layer} objects --- each carrying hand hardness, primary grain form, grain size, and layer thickness --- alongside a separate \texttt{density\_profile} list of density observations, each recording a depth, thickness, and density value. The \texttt{stability\_tests} object contains result lists for each test type; the \texttt{ECT} list is the relevant one for this analysis, with each entry recording the failure depth, number of taps, and propagation result.

% State the three key data handling decisions made in this analysis:
%   1. Density alignment: only density measurements attributed to a stratigraphic layer
%      are used; stand-alone density profile entries that cannot be unambiguously aligned
%      to a layer boundary are excluded
%   2. Grain form resolution

Three data handling decisions govern which observations are included in the analysis.

\textit{Density alignment.} SnowPilot allows density to be entered either as an attribute of an individual stratigraphic layer or as a separate density profile. When entered as a separate profile, a density observation cannot always be unambiguously matched to a stratigraphic layer boundary because the recorded depth and thickness intervals do not always coincide with layer boundaries. This analysis uses only density observations whose recorded depth and thickness match a layer boundary exactly; density profile entries that cannot be matched are excluded. Of the 19{,}877 density observations in the database, 10{,}692 (2.8\,\% of all layers) could be unambiguously aligned and are carried forward as measured density on the corresponding layer (see also Table~\ref{tab:snowpilot}, footnote~b).

\textit{Grain form resolution.} SnowPilot records grain forms using CAAML grain form codes, which may be either two-character basic grain class codes (e.g., \texttt{RG} for rounded grains, \texttt{FC} for faceted crystals) or longer sub-grain class codes (e.g., \texttt{RGmx} for mixed rounded grains, \texttt{FCxr} for depth-hoar-like facets). Each parameterization method defines a set of valid grain form codes. When a layer's recorded grain form code is not directly valid for a given method, the code is truncated to its two-character basic grain class prefix and re-evaluated; if the basic class is valid for the method, it is used. If neither the full code nor the truncated prefix is valid for the method, the method is inapplicable to that layer and the pathway fails silently for it.

\textit{Primary grain form.} Each layer in SnowPilot may carry both a primary and a secondary grain form. This analysis uses only the primary grain form for all parameterization calculations; secondary grain form information is not used.


\subsection{Measurement uncertainties}
\label{sec:uncertainties}

% SnowPilot stores nominal (point) values; real measurements carry uncertainty from
% both instrument/technique precision and spatial variability within the pit.

% Distinguish two uncertainty components explicitly:
%   1. Input measurement uncertainty: uncertainty in the field observation itself
%      (hand hardness, grain size, layer thickness, directly measured density, slope angle)
%   2. Method uncertainty: regression standard error from the original publication

% State that this paper reports input measurement uncertainty only:
%   Not all implemented methods provide regression errors, so direct comparison of
%   method uncertainty across parameterizations is not possible; reporting input
%   measurement uncertainty alone enables consistent comparison.

% Uncertainties are represented using the Python uncertainties package
% (lebigot2010uncertainties) as ufloat objects, enabling automatic first-order
% Taylor expansion uncertainty propagation through all calculations.

% Present adopted measurement uncertainties with justification:
%   - Hand hardness:    +/-0.67 HHI (+/-2 sub-classes)    [holler2010; author judgement]
%   - Slope angle:      +/-2 degrees                       [adapted from mccammonSLOPEMEASUREMENTHUMANS2023]
%   - Grain size:       +/-0.5 mm                          [author estimate]
%   - Layer thickness:  +/-5% relative                     [author estimate]
%   - Density (direct): +/-10% relative                    [conger2009, proksch2016]

% Note: uncertainty values were reviewed by an experienced field observer (birkeland2025).


\subsection{Parameterizations for mechanical properties}
\label{sec:parameterizations}

% Overview paragraph (write this paragraph first):
%   - Scope: we implement a selection of existing published parameterizations; we do not
%     develop new parameterizations or fit new regressions to the SnowPilot data
%   - Selection criteria: methods are included if they (a) take standard snowpit
%     observations as inputs and (b) represent well-cited recent work spanning a range
%     of applicable grain forms; exhaustive inclusion of every published attempt is
%     out of scope
%   - All methods return effective, isotropic elastic properties, consistent with the
%     isotropic assumption required by the slab models introduced in Sect. 1.1

\subsubsection{Density}

% Introduce density parameterizations:
%   - data_flow: direct measurement (coverage: ~2.8% of layers)
%   - geldsetzer: Geldsetzer & Jamieson (2000) -- hand hardness + grain form
%   - kim_jamieson_table2: Kim & Jamieson (2014) Table 2 -- hand hardness + grain form
%   - kim_jamieson_table5: Kim & Jamieson (2014) Table 5 -- hand hardness + grain form + grain size

% Note excluded grain forms: MM, SH, IF, and wet/melt forms except MFcr.

% [Table 2 near here: grain forms and hand hardness ranges covered by each density method,
%  comparing Geldsetzer & Jamieson (2000) and Kim & Jamieson (2014)]

\FloatBarrier

\subsubsection{Effective, isotropic Young's modulus}

% Note that all E-mod methods take density as input; density method choice propagates
% into E-mod estimation (method chaining). This coupling is formalized in Sect. 2.4.

% Present four methods:
%   - bergfeld:  Bergfeld et al. (2023) -- density + grain form (PP, DF, RG)
%   - kochle:    Kochle & Schneebeli (2014) -- density + grain form (RG, FC, DH)
%   - wautier:   Wautier et al. (2015) -- density + grain form (wide range), power law
%   - schottner: Schottner et al. (2026) -- density + grain form [confirm details from elastic_modulus.py]

% [Table 3 near here: applicability matrix -- grain form x E-mod method, with valid
%  density ranges; add Schottner column to draft in working_draft.tex]

\FloatBarrier

\subsubsection{Effective, isotropic Poisson's ratio}

% Present two methods:
%   - kochle:     Kochle & Schneebeli (2014) -- grain form only (RG, FC, DH)
%   - srivastava: Srivastava et al. (2016) -- grain form + density range (validity limits apply)

% Note that Wautier et al. (2015) calculated orthotropic nu but did not provide an
% effective isotropic parameterization; not implemented.

% Note that srivastava depends on density, creating a structural link between the
% density method choice and the nu method choice. This is explained in Sect. 2.4.1.

\subsubsection{Shear modulus}

% G is derived from E and nu via the standard isotropic relation G = E / (2(1+nu));
% no independent parameterization is implemented.
% Wautier et al. (2015) provides a direct G parameterization (future extension,
% noted in Sect. 4.5).

\subsubsection{Plate theory parameters for composite slabs}

% Introduce classical laminate theory applied to layered snow slabs
% (Weissgraeber & Rosendahl 2023).

% Assumptions: layers are homogeneous, isotropic, linearly elastic, and perfectly bonded.
% Define reduced stiffness per layer: Q_i = E_i / (1 - nu_i^2)

% Present equations for A11, B11, D11, A55 and their physical meanings:
%   A11 (N/mm): extensional stiffness         = sum Q_i * h_i
%   B11 (N):    bending-extension coupling    = (1/2) sum Q_i * (z_{i+1}^2 - z_i^2)
%   D11 (N*mm): bending stiffness             = (1/3) sum Q_i * (z_{i+1}^3 - z_i^3)
%   A55 (N/mm): transverse shear stiffness    = kappa * sum G_i * h_i, kappa = 5/6

% D11 is the headline parameter: bending stiffness governs the elastic energy in the
% slab available to drive crack propagation. It is the primary result of this analysis.


\subsection{Parameterization graph and calculation pathways}
\label{sec:graph}

% Opening connector: the parameterizations reviewed in Sect. 2.3 are implemented
% within a graph-based computational framework that systematically enumerates and
% executes all valid calculation pathways from snowpit observations to target mechanical
% parameters. This section describes the structure and operation of that framework.

The parameterizations described in Sect.~\ref{sec:parameterizations} are implemented within a directed acyclic graph (DAG) that encodes all valid calculation pathways from snowpit observations to target mechanical parameters. This section describes the structure of the graph, the algorithm for enumerating pathways, and the execution engine that applies each pathway to a slab.

\subsubsection{Representing parameterizations as a directed graph}

% Represent the full set of parameterizations as a directed acyclic graph (DAG):
%   - Parameter nodes: represent quantities -- either directly observed inputs
%     (hand_hardness, grain_form, grain_size, measured_density, layer_thickness) or
%     derived quantities (density, E, nu, shear_modulus, D11, etc.)
%     Parameter nodes implement OR logic: any one incoming method is sufficient to
%     compute the parameter.
%   - Merge nodes: represent a specific calculation method; implement AND logic --
%     all required inputs must be available before the method can be applied.
%     Each merge node has exactly one outgoing edge to its output parameter node.
%   - Edges: directed data flow (parameter to merge node) or method output
%     (merge node to output parameter node)

The parameterization graph contains two types of nodes. \textit{Parameter nodes} represent quantities --- either directly observed field measurements (hand hardness, grain form, grain size, directly measured density, and layer thickness) or derived parameters (density, Young's modulus, Poisson's ratio, shear modulus, and the plate-theory stiffness parameters A11, B11, D11, and A55). A parameter node implements OR logic: any single incoming method edge is sufficient to determine the parameter, so the graph represents all alternative estimation methods for a given quantity without requiring any one of them. \textit{Merge nodes} implement AND logic: each merge node represents a specific combination of inputs required by one or more calculation methods, and all inputs feeding a merge node must be available before the corresponding method can be applied. Each merge node has exactly one outgoing edge to its output parameter node.

Edges are directed and of two kinds. Data-flow edges carry observed or previously computed values forward from a parameter node to a merge node. Method edges carry the result of a calculation from a merge node to the output parameter node, and are labelled with the method name (for example, \textit{geldsetzer}, \textit{wautier}, or \textit{kochle}). The root node, representing the snowpit observation, connects by data-flow edges to all directly measured quantities.

A key structural feature of the graph is the shared \textit{density} node. The density node feeds both the Young's modulus methods and the Srivastava Poisson's ratio method through the same merge node. Because it is a single shared node, any pathway that uses the Srivastava method for Poisson's ratio must use the same density method as is used for Young's modulus --- there is no independent density choice for Poisson's ratio. This coupling, together with the four available density methods, four Young's modulus methods, and two Poisson's ratio methods, constrains D11 to $4 \times 4 \times 2 = 32$ unique calculation pathways. The complete graph structure is illustrated in Fig.~\ref{fig:dag}.

\subsubsection{Finding calculation pathways}

% Describe the pathway enumeration algorithm (find_parameterizations):
%   - Backward traversal: starts at the target parameter node, recurses toward
%     the snow_pit root node
%   - OR logic at parameter nodes: enumerate all incoming method edges independently;
%     each edge represents an alternative pathway for computing that parameter
%   - AND logic at merge nodes: require all input branches; Cartesian product of
%     input sub-pathways produces all valid combinations
%   - Dynamic programming: results for each node are memoized to avoid recomputing
%     shared subgraphs
%   - Deduplication: structurally different traversals that resolve to identical
%     (parameter, method) assignments are collapsed using a method fingerprint

% Result: find_parameterizations(graph, D11) returns 32 unique pathways.

All valid calculation pathways from the snowpit root to a target parameter are enumerated by a backward traversal algorithm. Beginning at the target parameter node, the algorithm recurses toward the root, applying the logic appropriate to each node type. At parameter nodes (OR logic), each incoming method edge is explored independently; each edge represents an alternative pathway for computing the parameter. At merge nodes (AND logic), sub-pathways from all required input branches are combined using a Cartesian product, so every combination of upstream method choices is represented in the output. Results for each node are memoized after their first computation, preventing redundant traversal of shared subgraphs when the same node appears in multiple branches of the DAG.

Following traversal, the resulting path trees are converted to a structured list of pathway objects and deduplicated. The deduplication step identifies structurally distinct traversals that resolve to identical sets of (parameter, method) assignments --- an artefact of the shared density node, which the Cartesian-product expansion may enumerate more than once --- and retains only the first occurrence of each unique method fingerprint. Applied to D11, the enumeration returns 32 unique pathways.

\subsubsection{Executing parameterizations}

% Describe the execution engine (ExecutionEngine):
%   - Input: a Slab object (ordered list of Layer objects) and a target parameter
%   - For each parameterization, execute in two phases:
%       Phase 1 (layer-level): compute density, E, nu, and G for each layer in sequence
%       Phase 2 (slab-level): compute A11, B11, D11, A55 using all layer properties
%   - Density caching: density is cached per layer for pathways sharing the same
%     density method, avoiding redundant computation within a slab
%   - Failure handling: a pathway fails silently for a given slab or layer if any
%     required input is missing or falls outside a method's validity range
%     (e.g., grain form not supported, density out of range); failures are recorded
%     with the reason
%   - Provenance tracking: the method used to compute each parameter is recorded
%     for every successful pathway--slab combination

For each enumerated pathway, the execution engine applies a two-phase procedure to a \texttt{Slab} object --- an ordered list of \texttt{Layer} objects carrying the measured field observations.

In the first phase (layer-level), the engine iterates over all layers in the slab and computes, in dependency order, the density, Young's modulus, Poisson's ratio, and shear modulus of each layer using the methods specified by the current pathway. Density values are cached per layer, keyed by layer index and method name, and reused across pathways that share the same density method, avoiding redundant computation. Parameters that depend on density (Young's modulus, Poisson's ratio via the Srivastava method, and shear modulus) are always computed afresh for each pathway, because their values --- including the propagated uncertainty budget --- differ between pathways that use different upstream density methods; caching them would silently collapse distinct uncertainty budgets into a single incorrect result.

In the second phase (slab-level), the engine assembles the per-layer elastic properties computed in the first phase and applies the classical laminate theory formulae (Sect.~\ref{sec:parameterizations}) to compute A11, B11, D11, and A55 for the slab as a whole.

A pathway fails silently for a given slab if any required field observation is absent, or if any method's validity conditions are not satisfied (for example, a grain form outside the method's applicable range, or a density value outside its regression validity interval). Failures are recorded with the reason, enabling post-hoc diagnosis of coverage losses. For successful pathway--slab combinations, the method used to compute each parameter on each layer is recorded for provenance tracing.


\subsection{Constructing slabs from ECT observations}
\label{sec:ectp_slabs}

% Scientific motivation: to demonstrate the framework at population scale, a slab
% definition is needed for each observation. ECTP (propagating Extended Column Test)
% results are used because they are the largest subset of SnowPilot pits that include
% a defined weak layer depth, providing a natural slab definition -- the layers above
% the failure depth. No claim is made about the relationship between estimated D11
% and ECT outcome; ECTP slabs serve purely as a source of field-observed slab
% stratigraphies for framework application.

To demonstrate the framework at population scale, a slab definition is required for each observation. Extended Column Tests with propagating fracture (ECTP) provide a natural slab definition: the recorded failure depth identifies the weak layer, and all snow above it constitutes a candidate slab. The ECTP subset is also the largest category of SnowPilot stability tests that includes a well-defined failure depth (Table~\ref{tab:snowpilot}). We emphasise that no claim is made about the relationship between estimated D11 and ECT outcome; ECTP results serve here solely as a large, community-sourced source of field-observed slab stratigraphies with identifiable weak layers.

% Describe the procedure:
%   1. Parse SnowPilot CAAML files into SnowPylot SnowPit objects using SnowPylot
%   2. For each ECTP result in a pit:
%      a. Identify the weak layer: match failure depth to layer depth_top and depth_bottom
%      b. Define the slab as all layers above the weak layer (not including the weak layer)
%      c. Create a SnowPyt-MechParams Slab object from those layers, carrying
%         measured hand hardness, grain form, grain size, thickness, and slope angle
%   3. A single pit can yield multiple slabs if it has multiple ECTP results

SnowPilot CAAML XML files were parsed into \texttt{SnowPit} objects using the SnowPylot library (Sect.~\ref{sec:snowpylot}). For each ECTP result in a pit, the weak layer was identified by matching the recorded failure depth to the depth boundaries of each stratigraphic layer: the weak layer is the layer whose upper boundary is at or shallower than the failure depth and whose lower boundary lies strictly below it. The slab was then defined as all layers above the weak layer (excluding the weak layer itself), and a \texttt{Slab} object was constructed from those layers, carrying the measured hand hardness, grain form, grain size, and thickness of each layer, and the pit slope angle. A single pit contributes more than one slab if it contains multiple ECTP results at different depths.

% Report resulting dataset size: [N] ECTP slabs from [N] unique pits.
% Note any exclusion criteria (e.g., slabs with zero qualifying layers).

This procedure yielded 14{,}951 ECTP slab stratigraphies. ECTP results for which no stratigraphic layers were identified above the failure depth were excluded. These 14{,}951 slabs form the population-scale dataset for the coverage and parameterization analyses presented in Sect.~\ref{sec:coverage}.


\section{Results}

\subsection{Coverage and data loss by parameterization}
\label{sec:coverage}

% STRUCTURE: lead with the punchline -- coverage is low -- then explain why.
% This section sets up the reader to interpret the per-parameter results
% in Sects. 3.2--3.5 with the coverage constraint in mind.
%
% Opening statement (the headline):
%   Even the highest-coverage pathways to D11 succeed for only ~5% of ECTP slabs.
%   State this number in the first sentence of the section.
%
% Explain the compounding mechanism:
%   Layer-level coverage for individual methods is substantially higher (~28--63%
%   for hand-hardness-based density methods), but slab-level coverage requires
%   ALL layers in a slab to succeed for a given pathway. Each additional required
%   field (grain size, density measurement) and each grain-form restriction compounds
%   the data loss. Illustrate with the data_flow pathways: ~2.8% of layers have a
%   direct density measurement, but slab-level coverage for data_flow pathways
%   is effectively zero because whole-slab coverage requires every layer to have
%   a density measurement.
%
% Present a summary figure or table: layer-level and slab-level coverage for each
% density, E-mod, and nu method, and for D11 across all 32 pathways.
% [Fig. 3 (coverage overview) belongs here]
%
% Framing for the reader: this coverage landscape is the first key finding --
% it establishes that the practical reach of field-data-driven mechanical models
% is narrow, regardless of which parameterization is chosen. Per-parameter detail
% follows in Sects. 3.2--3.5.

Even under the most favourable combination of parameterization methods, valid D11 estimates can be obtained for only 5.1\,\% of the 14{,}951 ECTP slab stratigraphies in the dataset. This coverage ceiling is the primary finding of the analysis and sets an upper bound on the practical reach of field-data-driven mechanical stability models applied to standard snowpit observations.

The low ceiling arises from a compounding of validity failures across each step of the calculation chain. At the layer level, individual methods cover substantial fractions of the 377{,}729 layers: hand-hardness-based density methods apply to 28--63\,\% of layers (Sect.~\ref{sec:density_results}), E-mod methods to 5--55\,\% depending on grain form acceptance and density range (Sect.~\ref{sec:emod_results}), and Poisson's ratio methods to 28--42\,\% (Sect.~\ref{sec:nu_results}). Slab-level coverage is systematically lower because a slab succeeds only if every one of its constituent layers passes all validity checks simultaneously for a given pathway. This multiplicative requirement transforms modest layer-level failure rates into severe slab-level failure rates. The measured-density (\texttt{data\_flow}) pathway illustrates the compounding starkly: 2.8\,\% of layers carry a density observation coinciding with a layer boundary, yet slab-level density coverage reaches only 0.7\,\% (109 slabs) because density must be present for every layer in the slab simultaneously. By the time the D11 calculation chain is complete, \texttt{data\_flow} pathways succeed for at most 17 slabs (0.1\,\%), as the requirement for a complete per-layer density measurement stack is nearly never met in practice. The complete coverage landscape across all density, E-mod, $\nu$, and D11 pathway combinations is shown in Fig.~\ref{fig:coverage}. Per-parameter detail—coverage fractions, method-to-method differences in applicable grain forms and density ranges, and the within-method spread in estimated values—is presented in Sects.~\ref{sec:density_results}--\ref{sec:d11_results}.

\FloatBarrier

\subsection{Density estimates across methods}
\label{sec:density_results}

% Layer-level analysis:
%   - Distributions of estimated density for each method (geldsetzer, kim_table2,
%     kim_table5) across all applicable layers
%   - Within-method uncertainty from input measurement propagation
%   - Method-to-method differences in estimated values (means differ by ~208--283 kg/m3)

At the layer level, the three empirical density methods differ in both coverage and the distribution of estimates (Fig.~\ref{fig:density_distributions}). \citet{kimjamieson2009} Table~2 achieved the broadest reach, yielding density estimates for 239{,}080 of 377{,}729 layers (63.3\,\%), with a mean of 208\,\si{kg.m^{-3}} and a mean relative uncertainty of 16.4\,\%. The \citet{geldsetzer2009} parameterization covered 203{,}677 layers (53.9\,\%), with a lower mean of 193\,\si{kg.m^{-3}} and a relative uncertainty of 17.1\,\%. \citet{kimjamieson2009} Table~5 was the most restrictive, applying to 106{,}370 layers (28.2\,\%), with a mean of 199\,\si{kg.m^{-3}} and the highest relative uncertainty at 18.4\,\%. The measured-density pathway (\texttt{data\_flow}) was applicable to only 10{,}692 layers (2.8\,\%)—those with a measured density observation coinciding with a layer boundary—but yielded a substantially higher mean of 283\,\si{kg.m^{-3}} and the lowest relative uncertainty of 10.0\,\%, reflecting that the remaining uncertainty arises from the $\pm$10\,\% density measurement protocol rather than from regression scatter. The higher mean density of the \texttt{data\_flow} subset likely reflects a sampling bias toward denser, harder layers where snow scientists chose to record density observations.

% Slab-level analysis:
%   - Coverage at slab level: fraction of ECTP slabs where all layers have a
%     calculable density for each method
%   - Any systematic differences in which slabs succeed across density methods

At the slab level, complete density coverage—meaning that every layer above the ECTP failure depth can be assigned a density estimate—is achievable for a smaller fraction of slabs than layer-level coverage alone would suggest, because a slab succeeds only if all of its constituent layers pass the grain form and hardness validity checks simultaneously. Of 14{,}951 ECTP slabs, \citeauthor{kimjamieson2009} Table~2 yields complete slab density estimates for 6{,}030 slabs (40.3\,\%), followed by \citet{geldsetzer2009} at 4{,}598 slabs (30.8\,\%), \citeauthor{kimjamieson2009} Table~5 at 1{,}170 slabs (7.8\,\%), and the measured-density pathway at 109 slabs (0.7\,\%). The ordering of methods by slab-level coverage mirrors their layer-level ordering, suggesting that coverage loss at the slab level is driven primarily by the same grain form restrictions that limit layer-level applicability, compounded by the requirement that no layer in a multi-layer slab may be excluded.

% Close with implications for downstream: the spread in density estimates across
% methods propagates directly into Young's modulus estimation, since all implemented
% E-mod methods take density as input.

The $\sim$90\,\si{kg.m^{-3}} spread in mean density across methods (193--283\,\si{kg.m^{-3}}) is consequential for downstream calculations, because all implemented Young's modulus parameterizations take density as their primary input. Even before any E-mod method uncertainty is considered, this method-to-method density disagreement introduces a systematic source of inter-pathway variance. Section~\ref{sec:emod_results} quantifies how this density spread propagates into Young's modulus estimates.


\FloatBarrier

\subsection{Young's modulus estimates across methods}
\label{sec:emod_results}

% Layer-level analysis across the 16 pathways (4 density x 4 E-mod methods):
%   - Distributions of estimated E for each E-mod method (aggregated across
%     applicable density sub-pathways, or for representative combinations)
%   - Large inter-method spread: wautier ~281--311 MPa vs. bergfeld/schottner ~7--19 MPa
%     (hand-hardness density sub-pathways); kochle intermediate at ~71--134 MPa
%   - Within-pathway relative uncertainty from input measurement propagation (23--92%)

The 16 calculation pathways to Young's modulus arise from combining 4 density sub-pathways with 4 E-mod methods (Fig.~\ref{fig:emod_distributions}). Each E-mod method imposes its own grain form validity requirements on top of the grain form restrictions already applied at the density step, so layer-level E-mod coverage is bounded above by the density coverage of the corresponding sub-pathway and generally falls below it. \citet{wautier2015} accepts the broadest set of grain form classes (DF, RG, FC, DH, MF) and a density range of 103--544\,\si{kg.m^{-3}}, resulting in the highest layer-level coverage among all E-mod methods: the \citeauthor{kimjamieson2009} Table~2 sub-pathway yields 208{,}331 of 377{,}729 layers (55.2\,\%). Its power-law exponent of $n = 2.34$ is the lowest among the four methods, so density uncertainty is amplified by a relatively modest factor; within-pathway relative uncertainties for \citeauthor{wautier2015} pathways are 33.8--38.5\,\% for hand-hardness density sub-pathways and 23.4\,\% for the measured-density sub-pathway. \citet{schottner2026} also achieves broad layer-level coverage (48.8\,\%) by accepting DF, RG, FC, DH, and SH grain forms, but its grain-form-specific coefficient pairs ($A$, $n$) carry large standard errors ($\sigma_A/A \approx 40$--$75\,\%$; $\sigma_n/n \approx 12$--$24\,\%$), and together with a high exponent ($n = 4.6$--$5.1$), produce within-pathway uncertainties of 77--91\,\%. \citet{kochle2014} accepts RG, DH, and MF grain forms over 150--450\,\si{kg.m^{-3}} and covers 19.0--27.5\,\% of layers for hand-hardness sub-pathways; its purely exponential form produces high sensitivity to density errors and within-pathway uncertainties of 67--92\,\%. \citet{bergfeld2023} is the most restrictive, accepting only PP, RG, and DF grain forms within 110--363\,\si{kg.m^{-3}}, covering 23.7--23.9\,\% of layers for the two highest-coverage hand-hardness sub-pathways, with within-pathway uncertainties of 72--77\,\%.

The inter-method spread in layer-level mean $E$ is large (Fig.~\ref{fig:emod_distributions}). For hand-hardness density sub-pathways, \citeauthor{wautier2015} produces the highest mean values (281--311\,\si{MPa}), followed by \citeauthor{kochle2014} (71--134\,\si{MPa}), with \citeauthor{schottner2026} and \citeauthor{bergfeld2023} yielding substantially lower means of 12--19\,\si{MPa} and 7--13\,\si{MPa}, respectively. The measured-density sub-pathway yields higher means across all E-mod methods (wautier 656\,\si{MPa}, kochle 250\,\si{MPa}, schottner 83\,\si{MPa}, bergfeld 27\,\si{MPa}), consistent with the elevated mean density of measured-density layers (Sect.~\ref{sec:density_results}).

% Slab-level analysis: coverage at slab level for each pathway combination.

At the slab level, complete E-mod coverage requires every layer in a slab to obtain a valid E-mod estimate in addition to a valid density estimate. The highest slab-level coverage is achieved by \citeauthor{schottner2026} pathways: both the \citeauthor{geldsetzer2009} $\to$ schottner and \citeauthor{kimjamieson2009} Table~2 $\to$ schottner pathways succeed for 2{,}814 of 14{,}951 ECTP slabs (18.8\,\%). The \citeauthor{kimjamieson2009} Table~2 $\to$ \citeauthor{wautier2015} pathway follows at 2{,}134 slabs (14.3\,\%), while \citeauthor{kochle2014} pathways cover only 41--530 slabs (0.3--3.5\,\%) and measured-density sub-pathways fewer than 60 slabs (0.1--0.4\,\%). The full 16-pathway coverage matrix is presented in Fig.~\ref{fig:coverage}.

% Close with implications for downstream: the large inter-method spread in E estimates
% is the dominant driver of D11 spread.

The roughly 15--40-fold spread in mean layer-level $E$ between \citeauthor{wautier2015} ($\sim$300\,\si{MPa}) and \citeauthor{bergfeld2023}/\citeauthor{schottner2026} ($\sim$8--15\,\si{MPa}) for hand-hardness density sub-pathways is the dominant source of variance in D11 across the 32 pathways. Because D11 is linear in the layer-by-layer $E$ values through the CLT integration, this inter-method E factor propagates directly as the same multiplicative factor in D11. Poisson's ratio, examined in Sect.~\ref{sec:nu_results}, varies less across methods and contributes secondarily to D11 spread (Sect.~\ref{sec:d11_results}).


\FloatBarrier

\subsection{Poisson's ratio estimates across methods}
\label{sec:nu_results}

% Layer-level analysis for kochle and srivastava:
%   - kochle: nu by grain form (RG: 0.171, FC: 0.130, DH: 0.087); no density dependence
%   - srivastava: nu by grain form group with density validity ranges; produces
%     4 sub-pathways (one per density method) -- report coverage and values for each
%   - Inter-method spread is smaller than for density or E-mod

Five calculation pathways to Poisson's ratio arise from combining two methods with their respective density sub-pathway requirements. The \citet{kochle2014} method requires no density input: it assigns a grain-form-specific mean $\nu$ value (RG: $0.171\pm0.026$; FC: $0.130\pm0.040$; DH: $0.087\pm0.063$) derived from the same FE simulations of m-CT snow microstructure used for the kochle E-mod method. Because density is not involved, the kochle $\nu$ pathway generates a single branch and achieves the broadest layer-level coverage: 159{,}120 of 377{,}729 layers (42.1\,\%), with an overall mean $\nu$ of 0.144 (reflecting the grain-form composition of the RG-/FC-/DH-applicable subset). The \citet{srivastava2016} method also uses a grain-form-specific lookup (RG: $0.191\pm0.008$; PP and DF: $0.132\pm0.053$; FC and DH: $0.170\pm0.020$), but accepts five grain form classes and imposes a density validity threshold ($\rho > 200$\,\si{kg.m^{-3}}). Because the density threshold is enforced by the graph, the srivastava pathway requires a density sub-pathway, producing four branches—one per density estimation method. Layer-level coverage for the srivastava hand-hardness sub-pathways is 28.0--29.4\,\% (105{,}757--110{,}937 layers), reflecting both grain form and density-threshold restrictions; the \citeauthor{kimjamieson2009} Table~5 and measured-density sub-pathways cover 10.8\,\% and 1.5\,\%, respectively. Mean $\nu$ across srivastava sub-pathways is 0.170--0.178. A notable property of both methods is that propagated uncertainty (with method uncertainty excluded) is exactly zero: because $\nu$ is assigned as a fixed value per grain form category rather than computed as a function of continuous density, the density measurement uncertainty does not propagate into $\nu$.

% Note that nu contributes less to D11 uncertainty than E-mod method choice, but the
% srivastava--kochle difference still contributes to the overall D11 spread (Sect. 3.5).

At the slab level, kochle $\nu$ succeeds for 944 of 14{,}951 ECTP slabs (6.3\,%), the highest slab-level coverage among all $\nu$ pathways and higher than most E-mod pathway–slab combinations. Srivastava slab-level coverage reaches 3.1\,\% for the \citeauthor{geldsetzer2009} and \citeauthor{kimjamieson2009} Table~2 sub-pathways (461 slabs each), 0.3\,\% for Table~5, and 0.1\,\% for measured density. The inter-method spread in $\nu$ ($\sim$0.144 versus $\sim$0.170--0.178, a difference of $\sim$24\,\%) is smaller than the inter-method spread in $E$ (up to 40-fold) and contributes correspondingly less to D11 variance. However, the kochle–srivastava $\nu$ contrast does introduce a systematic offset in D11 that adds to the E-mod-driven spread examined in Sect.~\ref{sec:d11_results}.

% Close with implications for downstream: Poisson's ratio method choice is not
% independent of density method choice when srivastava is used, because both
% draw from the same shared density node in the graph (Sect. 2.4.1). This structural
% coupling is reflected in the 32-pathway count for D11.

The structural coupling introduced by the shared density node (Sect.~\ref{sec:dag}) has a direct consequence for D11 pathway enumeration. When srivastava $\nu$ is selected, the density sub-pathway is already determined by the E-mod step (both share the same \texttt{merge\_density\_grain\_form} node in the DAG), so no additional density branching occurs. When kochle $\nu$ is selected, $\nu$ is independent of the density sub-pathway entirely. In either case, the $\nu$ method choice adds exactly one branch to each of the 16 E-mod pathways, yielding $16 \times 2 = 32$ unique D11 pathways in total. This means that srivastava does not multiply the pathway count by 4 (one per density sub-pathway) at the D11 level—the coupling collapses what would otherwise be additional branching into a single constrained choice per E-mod pathway.

\FloatBarrier

\subsection{Bending stiffness (D11) of ECT slabs across all pathways}
\label{sec:D11_results}

% STRUCTURE: three findings presented in order of impact.
% Frame each paragraph as a finding, not just a report.
%
% Paragraph 1 -- FINDING 1: Standard field data are insufficient to compute D11
%   for the vast majority of slabs.
%   - The highest-coverage pathways succeed for approximately 5% of slabs
%     (4 pathways at 756/14,951 = 5.1%: geldsetzer/kim_table2 × wautier/schottner × kochle)
%   - Direct-measurement (data_flow) pathways succeed for only ≤17 slabs (0.1%)
%   - 95% of ECTP slabs have no valid D11 estimate; this is an observation gap
%   [Fig. 7 D11 distributions belongs here]
%
% Paragraph 2 -- FINDING 2: For slabs where D11 can be computed, the result
%   depends strongly on method choice.
%   - Across the 4 highest-coverage pathways, D11 spans ~26× (8.7e7 to 2.3e9 N·m)
%     driven entirely by E-mod method choice (wautier vs schottner)
%   - Full 32-pathway range: ~3.6e7 to ~5.0e9 N·m
%   - E-mod method choice is the dominant contributor; ν method adds ~1.4×;
%     density method adds ~1.0–1.1×
%
% Paragraph 3 -- FINDING 3: Within-pathway uncertainty is also large and compounds
%   the inter-pathway spread.
%   - Within-pathway uncertainty: 23–97%; for highest-coverage pathways: 37–75%
%   - The 26× inter-pathway spread substantially exceeds within-pathway ±1σ (factor 1.4–2.7×)
%   - Together: D11 estimates carry substantial uncertainty from both method choice
%     and measurement imprecision

For the 14{,}951 ECTP slab stratigraphies assembled from the SnowPilot database, valid D11 estimates can be obtained for at most 5.1\,\% of slabs under any implemented pathway (Fig.~\ref{fig:d11_distributions}). The maximum of 756 slabs is achieved by four pathways sharing a common structural feature: the kochle Poisson's ratio method (which requires only grain form, not a density estimate), combined with either the \citeauthor{wautier2015} or \citeauthor{schottner2026} E-mod method (both accepting five grain form classes) and either the \citeauthor{geldsetzer2009} or \citeauthor{kimjamieson2009} Table~2 density sub-pathway (the two broadest hand-hardness methods). These combinations maximise coverage by avoiding the grain size requirement of \citeauthor{kimjamieson2009} Table~5, the density measurement requirement of the \texttt{data\_flow} sub-pathway, and the grain form restrictions of the \citeauthor{kochle2014} and \citeauthor{bergfeld2023} E-mod methods. Pathways that require a direct density observation (\texttt{data\_flow}) succeed for at most 17 slabs (0.1\,\%), because a layer-by-layer density measurement must be present for every layer in a slab simultaneously, and such complete density coverage within a single slab is rare in the SnowPilot database. The remaining 94.9\,\% of ECTP slabs—spanning more than a decade of observations from avalanche professionals and snow scientists—yield no valid D11 estimate under any of the 32 implemented pathways. This is not a limitation of the framework: the framework successfully identifies valid pathways wherever one exists. It is a gap in the observations: standard snowpit recording protocols do not routinely capture the combination of grain form, hand hardness, layer thickness, and Poisson's-ratio-compatible grain types that D11 estimation requires for every layer.

For the slabs where D11 can be computed, the result depends strongly on which parameterization pathway is used (Fig.~\ref{fig:d11_distributions}). Considering the four highest-coverage pathways (756 slabs each, all using kochle $\nu$), mean slab D11 ranges from $8.7 \times 10^7$\,\si{N.m} (\citeauthor{kimjamieson2009} Table~2 $\to$ \citeauthor{schottner2026} $\to$ kochle) to $2.3 \times 10^{9}$\,\si{N.m} (\citeauthor{geldsetzer2009} $\to$ \citeauthor{wautier2015} $\to$ kochle)—a factor of $\sim$26 arising solely from the choice of E-mod method applied to the same set of field observations. Across all 32 pathways, mean D11 spans from $\sim$3.6 $\times 10^7$\,\si{N.m} (data\_flow $\to$ \citeauthor{bergfeld2023} $\to$ kochle, 6 slabs) to $\sim$5.0 $\times 10^{9}$\,\si{N.m} (data\_flow $\to$ \citeauthor{wautier2015} $\to$ srivastava, 12 slabs). Isolating the contribution of each method choice: fixing density and $\nu$ method, the E-mod method contributes a 26-fold spread (wautier vs.~schottner or bergfeld); fixing density and E-mod method, the $\nu$ method (kochle vs.~srivastava) contributes a factor of $\sim$1.4; and fixing E-mod and $\nu$ method, density method choice contributes a factor of $\sim$1.0--1.1. E-mod parameterization choice is thus the dominant source of inter-pathway D11 variance, with $\nu$ and density method choices playing secondary roles.

Within any single pathway, D11 estimates also carry substantial uncertainty attributable to propagated field measurement imprecision. Within-pathway relative uncertainties (input measurement propagation only, excluding method regression uncertainty) range from 23\,\% (data\_flow $\to$ \citeauthor{wautier2015} $\to$ kochle, 17 slabs) to 97\,\% (geldsetzer $\to$ kochle $\to$ srivastava, 344 slabs); for the four highest-coverage pathways, within-pathway uncertainties are 37--75\,\%. At $\pm$1\,$\sigma$, this within-pathway uncertainty represents a multiplicative factor of 1.4--2.7 on a single D11 estimate, substantially smaller than the 26-fold inter-pathway spread but not negligible. The inter-pathway spread (26-fold) and the within-pathway spread ($\pm$1\,$\sigma$ factor 1.4--2.7) are therefore distinct and complementary components of the total D11 uncertainty budget: the former reflects disagreement among published parameterizations about the physical relationship between field observables and mechanical properties, while the latter reflects the limited precision of those field observables themselves. Any D11 estimate derived from field observations carries both sources simultaneously, and neither can be reported without the other if the estimate is to be used as input to a quantitative mechanical stability model.

\FloatBarrier


\section{Discussion}

\subsection{Implications for field-data-driven mechanical stability models}
\label{sec:discussion_implications}

% STRUCTURE: this is the central discussion section -- the place where the findings
% connect to the operational question. Write it as an argument, not a report.
% The three findings from §3.5 each have an implication; state them explicitly.
%
% Implication of FINDING 1 (coverage):
%   A mechanical stability model that requires D11 as input can be applied to field
%   data from roughly 1 in 20 standard snowpit observations, at best. This is not
%   a limitation of any particular model -- it is a limitation of what is routinely
%   recorded. Practitioners who wish to use these models in the field will find that
%   standard observation protocols are generally insufficient without additional
%   measurements (e.g., direct density on every layer, grain size throughout).
%   This raises the question of whether mechanical models, as currently parameterized,
%   are operationally viable for widespread field use.
%
% Implication of FINDING 2 (inter-pathway spread):
%   When D11 is computable, the result varies by a factor of 3--4 depending solely
%   on which published parameterization is used. For a model such as WEAC that uses
%   D11 directly as an input, this spread would propagate into stability estimates:
%   identical field observations, processed by two practitioners using different
%   (but individually defensible) method choices, would produce fundamentally different
%   outputs. This is not a failure of the methods -- it reflects genuine scientific
%   uncertainty about how to convert field observations to mechanical properties.
%   But it means that model outputs are implicitly method-dependent in a way that is
%   rarely acknowledged.
%   Note that this finding does not invalidate the mechanical models themselves --
%   it reveals and quantifies a gap in the estimation of their inputs at population
%   scale for the first time.
%
% Implication of FINDING 3 (within-pathway uncertainty):
%   Even when method choice is fixed, relative uncertainty from input measurement
%   propagation is 40--90%. Combined with the inter-pathway spread, both sources
%   of uncertainty must be quantified and reported before field-data-driven mechanical
%   model outputs can be used with appropriate confidence. Stability conclusions
%   derived without acknowledging these uncertainties should be interpreted with caution.
%
% Synthesis (close the loop to §1.1):
%   Together, the three findings establish that the inputs to field-data-driven
%   mechanical stability models are simultaneously (a) rarely achievable from standard
%   observations and (b) poorly constrained when achievable. This does not mean
%   mechanical models should be abandoned -- it means the field needs validated
%   parameterizations, better measurement protocols, and transparent uncertainty
%   reporting before these models can be used reliably for operational stability
%   assessment.

\subsection{Toward a preferred calculation pathway}
\label{sec:discussion_best_pathway}

% Discuss whether a single "best" pathway can be identified and how one might approach it:
%   - Coverage alone is insufficient: higher coverage does not imply higher accuracy
%   - Without ground truth for mechanical properties in field conditions, validation
%     of individual pathway accuracy is not currently possible from field data alone
%   - Possible selection criteria: alignment with controlled laboratory measurements,
%     physical plausibility of estimated ranges, internal consistency across parameters

\subsection{Should mechanical properties be measured directly in the snowpit?}
\label{sec:discussion_direct_measurement}

\subsection{Limitations and assumptions}
\label{sec:limitations}

% SnowPilot data quality:
%   - Community-sourced data spans a wide range of observer expertise and completeness
%   - Nominal (point) values recorded; spatial variability within a pit is not captured
%   - Observer bias in grain form assignment is possible and not quantifiable from
%     the database alone

% Methodological assumptions:
%   - Isotropic elastic behavior assumed at the layer level; real snow is anisotropic,
%     particularly for faceted crystals and depth hoar
%   - Perfect bonding between slab layers assumed in classical laminate theory
%   - Effective isotropic properties applied uniformly; validity of the isotropic
%     assumption varies by grain form

% Parameterization coverage gaps:
%   - Melt-freeze crusts (MFcr) are excluded from most E-mod and nu parameterizations
%   - Machine-made snow (MM), surface hoar (SH), and ice formations (IF) are not
%     covered by any implemented E-mod or nu method

\subsection{Extensions and future work}
\label{sec:extensions}

% Framework extensibility:
%   - Adding a new parameterization requires: (1) implement the calculation function,
%     (2) add a node or edge to the graph definition, (3) register the method in
%     the execution dispatcher
%   - The graph and execution engine do not require modification; only the
%     parameter-specific module and graph definition are touched
%   - This design means newly published parameterizations can be integrated and
%     compared against existing methods with minimal effort

% Priority extensions:
%   - Weak layer failure parameters (shear strength, critical energy release rate Gc):
%     would enable full WEAC calculation from field data; PST observations provide
%     the natural entry point (critical cut length -> Gc via Weissgraeber & Rosendahl 2023)
%   - Direct G parameterization: Wautier et al. (2015) provides a direct shear modulus
%     parameterization not currently implemented; would offer an alternative to the
%     derived G = E / (2(1+nu))
%   - Elastic properties for melt-freeze crusts: no suitable parameterization exists
%     to our knowledge; future laboratory work is recommended

% Research opportunities identified by this analysis:
%   - The 3--4x inter-pathway spread in D11 motivates controlled field campaigns that
%     measure E and D11 directly alongside standard snowpit observations, enabling
%     validation of individual parameterizations against measured values
%   - As mechanical models grow more complex (more parameters, more input requirements),
%     systematic frameworks become increasingly necessary to understand the combined
%     effect of parameter uncertainty on model predictions


\section{Conclusions}

% STRUCTURE: findings first, framework last. Each conclusion is a complete, standalone
% statement that a reader should be able to cite independently.
%
% 1. FINDING -- Coverage: Standard snowpit observations are insufficient to compute
%    D11 for the vast majority of field datasets. The highest-coverage calculation
%    pathway produces a valid D11 estimate for approximately 5% of 14,951 ECTP slabs
%    drawn from 10 years of SnowPilot data; pathways requiring direct density
%    measurement succeed for fewer than 0.1%. This low coverage reflects limitations
%    in what is routinely recorded during standard pit observations, not a failure
%    of any specific parameterization.
%
% 2. FINDING -- Inter-pathway spread: Where D11 can be computed, it varies by a factor
%    of 3--4 across the 32 valid pathways for the same slab, depending entirely on
%    method choice. Young's modulus parameterization is the dominant contributor to
%    this spread. Identical field observations processed with different (individually
%    defensible) methods yield fundamentally different D11 estimates.
%
% 3. FINDING -- Within-pathway uncertainty: Within-pathway relative uncertainty from
%    input measurement propagation alone is 40--90%. This compounds the inter-pathway
%    spread: method choice and measurement imprecision each contribute uncertainties
%    of comparable magnitude to the total D11 uncertainty budget.
%
% 4. IMPLICATION -- Risk of over-reliance: Together, low coverage and high uncertainty
%    from both sources establish that field-data-driven mechanical stability models
%    rest on inputs that are both rarely achievable from standard observations and
%    poorly constrained when achievable. Stability conclusions derived from these
%    model outputs should be accompanied by explicit uncertainty quantification that
%    accounts for both parameterization sensitivity and measurement imprecision.
%
% 5. TOOL -- SnowPyt-MechParams: The framework that produced these results is
%    openly available and explicitly extensible. For D11, it enumerates 32 unique
%    pathways (4 density × 4 E-mod × 2 nu), constrained by the shared density node
%    in the parameterization graph. Adding a new parameterization requires only
%    implementing the calculation function, adding an edge to the graph, and
%    registering the method in the dispatcher -- enabling future comparison as
%    new methods are published.


%% The following commands are for the statements about the availability of data sets and/or software code corresponding to the manuscript.

\codeavailability{SnowPyt-MechParams is openly available at [URL TBD]. SnowPylot is openly available at [URL TBD].}

\dataavailability{The SnowPilot database is publicly accessible at snowpilot.org. The CAAML export files used in this analysis cover the period 1 September 2014 to 31 August 2025.}

\authorcontribution{TEXT} %% this section is mandatory

\competinginterests{TEXT} %% this section is mandatory even if you declare that no competing interests are present

\begin{acknowledgements}
TEXT
\end{acknowledgements}

\financialsupport{TEXT} %% list all funding sources here


%% REFERENCES
\bibliographystyle{copernicus}
\bibliography{references}


\end{document}
